% Generated draft: requires mathematical and source review.
\documentclass[11pt]{article}
\usepackage[margin=25mm]{geometry}
\usepackage{fontspec}
\setmainfont{Latin Modern Roman}
\usepackage{amsmath,amssymb,mathtools,unicode-math}
\setmathfont{Latin Modern Math}
\usepackage{enumitem,needspace,xurl,hyperref}
\setlength{\parindent}{0pt}
\setlength{\parskip}{5pt}
\setlength{\emergencystretch}{3em}
\newcommand{\N}{\mathbb{N}}
\newcommand{\Z}{\mathbb{Z}}
\newcommand{\Q}{\mathbb{Q}}
\newcommand{\R}{\mathbb{R}}
\newcommand{\C}{\mathbb{C}}
\newcommand{\F}{\mathbb{F}}
\newcommand{\A}{\mathbb{A}}
\newcommand{\PP}{\mathbb{P}}
\newcommand{\E}{\mathbb{E}}
\newcommand{\eps}{\varepsilon}
\newcommand{\dd}{\,\mathrm{d}}
\begin{document}
\section*{UnsolvedMath ID: 30000396}
\textbf{Source problem code:} catalogue entry
\subsection{Definitions and mathematical statement}
The source asks for a moduli space of Lie algebra structures up to equivalence, with local structure at a fixed algebra governed by its deformations. It does not specify the coefficient field, fixed underlying vector space/dimension, finite versus infinite-dimensional category, allowed families or moduli universal property.

For a fixed finite-dimensional vector space $V$ over a field $k$, the conventional parameter space is the bracket locus
\[
 \operatorname{Lie}(V)=\{\mu\in\operatorname{Hom}_k(\Lambda^2V,V):
 \mu(x,\mu(y,z))+\mu(y,\mu(z,x))+\mu(z,\mu(x,y))=0\ \forall x,y,z\}.
\]
The base-change action is $(g\cdot\mu)(x,y)=g\mu(g^{-1}x,g^{-1}y)$, so its $\operatorname{GL}(V)$-orbits are isomorphism classes in this fixed setting. This illustrative locus and set quotient alone do not specify a fine, coarse, stack-theoretic, analytic or topological moduli object.

For a fixed bracket $\mu$, the marked local deformation problem assigns to a local Artinian $k$-algebra with residue field $k$ the brackets on $V$ extended to that algebra reducing to $\mu$, modulo basis changes reducing to the identity. Formal deformations over $k[[t]]$ are a corresponding complete-local version. A rigorous interpretation of ``local structure reflected by deformations'' would specify how this functor or groupoid is recovered from a formal neighborhood or a versal deformation base. Marked deformations, unmarked isomorphism classes and their automorphisms are distinct; a set quotient is not automatically a representing fine moduli space.

The question is to construct an appropriate moduli object with the stated local-to-deformation relation in the intended category, not merely to put an arbitrary topology on the set of classes. Infinite-dimensional or graded/topological Lie algebras would require different parameter and family conditions, which cannot be inferred from the finite-dimensional illustration.

The literature triage checked 21 August 2026 calls the extraction incomplete because the structures to be parameterized are omitted. The present supplied sentence does identify Lie algebra structures, but still leaves their category and moduli specification open. This discrepancy is retained, and no universal representability claim or independently verified open theorem is asserted.
\subsection{Short English statement}
Can equivalence classes of Lie brackets form a natural moduli object whose local neighborhoods encode deformations of each algebra? A precise question must fix the Lie-algebra category, families, equivalences and meaning of local moduli. The supplied assessment still regards the scope as incomplete.
\subsection{Sources}
\begin{itemize}
\item Local UnsolvedMath registry, permanent ID 30000396, exact source code catalogue entry; \emph{Moduli Spaces of Lie Algebra Structures}. Proposed year 2006, status open. The verification reports normalization preserving source-level content; no direct report inspection is claimed here.
\item \emph{Deformations and Contractions in Mathematics and Physics (2006)}, Oberwolfach Reports, DOI 10.4171/owr/2006/03, pp. 119--186, \url{https://doi.org/10.4171/owr/2006/03}.
\item Supplied literature triage checked 2026-08-21, citing only ``original Oberwolfach report'', \url{https://doi.org/10.4171/owr/2006/03}. It says the structures to be parameterized are omitted, while this retained sentence names Lie algebra structures; it provides no precise resolution. The dated OWR-derived curation requests independent verification of time-sensitive status and cited claims.
\end{itemize}
\end{document}
